\par\Needspace{6\baselineskip}\medskip
\SourceLeaf{142}\pdfbookmark[2]{102.}{article-102}
\begin{center}102.\end{center}
Quum in hyperbola perinde ut in ellipsi $y$ necessario esse debeat positiva, solutio aequationis 16 hic quoque ambiguitati obnoxia esse nequit\footnote{Vix opus erit monere, tabulam nostram II in hyperbola perinde ut in ellipsi ad solutionem huius aequationis adhiberi posse, quamdiu $h$ ipsius limites non egrediatur.}: sed respectu aequationis $16^*$ hic paullo aliter ratiocinandum est quam in ellipsi. Ex aequationum theoria facile demonstratur, pro valore positivo ipsius $H$\HistoricalNoteDoubleUnbracketed{Quantitas $H$ manifesto fieri nequit negativa, nisi fuerit $\zeta>\tfrac16$; tali autem valori ipsius $\zeta$ responderet valor ipsius $z$ maior quam $2{,}684$, adeoque limites huius methodi longe egrediens.} hanc aequationem (siquidem ullam radicem realem positivam habeat) cum una radice negativa duas positivas habere, quae vel ambae aequales erunt puta $=\tfrac16\BareRoot5-\tfrac16=0{,}20601$, vel altera hoc limite maior altera minor. Iam in problemate nostro (suppositioni superstructo, $z$ esse quantitatem haud magnam, saltem non maiorem quam $0{,}3$, ne tabulae tertiae usu destituamur) necessario semper radicem maiorem accipiendam esse sequenti modo demonstramus. Si in aequatione $13^*$ pro $M$ substituitur $Y\BareRoot{(L+z)}$, prodit
\[
Y+1=(L+z)Z>(1+z)Z,
\]
\noindent sive
\[
Y>\tfrac13-\frac4{3\mathpunct{.}5}z+\frac{4\mathpunct{.}6}{3\mathpunct{.}5\mathpunct{.}7}zz-\frac{4\mathpunct{.}6\mathpunct{.}8}{3\mathpunct{.}5\mathpunct{.}7\mathpunct{.}9}z^3+\text{etc.},
\]
\noindent unde facile concluditur, pro valoribus tam parvis ipsius $z$, quales hic supponimus, semper fieri debere $Y>0{,}20601$. Revera calculo facto invenimus, ut $(1+z)Z$ huic limiti aequalis fiat, esse debere $z=0{,}79858$: multum vero abest, quin methodum nostram ad tantos valores ipsius $z$ extendere velimus.
